%% Reusable TikZ pieces for the illustrated figures (robots, token bars, coins,
%% cards, speech bubbles, step headers). Read from macros.tex. Macros only:
%% nothing here changes the page layout.
%% Every piece takes an optional scale [s] and a position {(x,y)}.
\usetikzlibrary{arrows.meta,shapes.geometric,decorations.pathmorphing,calc}

\definecolor{sggold}{HTML}{D9A030}
\definecolor{sgoff}{HTML}{9A9A9A}

\tikzset{
  sgfont/.code={\pgfmathsetmacro\sgfs{#1*1.2}%
    \edef\sgtmp{\noexpand\tikzset{font=\noexpand\sffamily\noexpand\fontsize{#1}{\sgfs}\noexpand\selectfont}}\sgtmp},
  sgfont/.default=7,
  sgthin/.style={line width=.45pt},
  sgflow/.style={-{Stealth[length=4pt,width=3.4pt]},line width=.9pt},
}

%% \sgrobot[s]{(x,y)}{colour}{label}: friendly robot, origin = head centre.
%% The label sits above the antenna; the body ends about 0.45 below the origin.
\newcommand{\sgrobot}[4][1]{%
  \begin{scope}[shift={#2},scale=#1]
    \draw[#3,line width=.6pt] (0,.19)--(0,.30);
    \fill[#3] (0,.33) circle (.04);
    \draw[#3,line width=.6pt] (-.15,-.27)--(-.27,-.37);
    \draw[#3,line width=.6pt] (.15,-.27)--(.27,-.37);
    \fill[#3] (-.28,-.38) circle (.03) (.28,-.38) circle (.03);
    \filldraw[draw=#3,fill=#3!14,line width=.6pt,rounded corners=1.2pt]
      (-.15,-.20) rectangle (.15,-.42);
    \draw[#3!60,line width=.5pt] (-.06,-.29)--(.06,-.29);
    \filldraw[draw=#3,fill=#3!14,line width=.65pt,rounded corners=2.2pt]
      (-.25,-.19) rectangle (.25,.19);
    \filldraw[draw=#3!80,fill=white,line width=.4pt,rounded corners=1.5pt]
      (-.18,-.12) rectangle (.18,.12);
    \fill[#3!75!black] (-.075,.02) circle (.034) (.075,.02) circle (.034);
    \draw[#3!75!black,line width=.5pt] (-.04,-.055) .. controls (0,-.08) .. (.04,-.055);
  \end{scope}%
  \node[text=#3!85!black,font=\sffamily\bfseries\fontsize{8}{8}\selectfont,
    inner sep=0pt,anchor=south] at ($#2+(0,#1*.39)$) {#4};%
}

%% \sgrobotoff[s]{(x,y)}{label}: shut-down robot (grey, closed eyes).
\newcommand{\sgrobotoff}[3][1]{%
  \begin{scope}[shift={#2},scale=#1]
    \draw[sgoff,line width=.6pt] (0,.19)--(0,.30);
    \draw[sgoff,line width=.5pt,fill=white] (0,.33) circle (.04);
    \draw[sgoff,line width=.6pt] (-.15,-.27)--(-.22,-.41);
    \draw[sgoff,line width=.6pt] (.15,-.27)--(.22,-.41);
    \filldraw[draw=sgoff,fill=black!8,line width=.6pt,rounded corners=1.2pt]
      (-.15,-.20) rectangle (.15,-.42);
    \filldraw[draw=sgoff,fill=black!8,line width=.65pt,rounded corners=2.2pt]
      (-.25,-.19) rectangle (.25,.19);
    \filldraw[draw=sgoff,fill=black!4,line width=.4pt,rounded corners=1.5pt]
      (-.18,-.12) rectangle (.18,.12);
    \draw[sgoff,line width=.6pt] (-.11,.02)--(-.04,.02) (.04,.02)--(.11,.02);
    \draw[sgoff,line width=.5pt] (-.04,-.06)--(.04,-.06);
  \end{scope}%
  \node[font=\sffamily\bfseries\fontsize{8}{8}\selectfont,text=sgoff,
    inner sep=0pt,anchor=south] at ($#2+(0,#1*.39)$) {#3};%
}

%% \sgbar[s]{(x,y)}{colour}{filled}{cells}: token bar centred at (x,y).
\newcommand{\sgbar}[5][1]{%
  \begin{scope}[shift={#2},scale=#1]
    \pgfmathsetmacro{\sgbarw}{#5*.115-.025}
    \foreach \sgi in {1,...,#5}{%
      \pgfmathsetmacro{\sgbx}{-\sgbarw/2+(\sgi-1)*.115}
      \ifnum\sgi>#4\relax
        \draw[#3!70,line width=.35pt,fill=white] (\sgbx,-.04) rectangle ++(.09,.08);
      \else
        \filldraw[draw=#3!80!black,fill=#3!80,line width=.35pt] (\sgbx,-.04) rectangle ++(.09,.08);
      \fi}
  \end{scope}%
}

%% \sgcoins[s]{(x,y)}{n}: stack of n coins, bottom coin centred at (x,y).
\newcommand{\sgcoins}[3][1]{%
  \begin{scope}[shift={#2},scale=#1]
    \foreach \sgi in {1,...,#3}{%
      \filldraw[draw=sggold!70!black,fill=sggold!55,line width=.4pt]
        (-.12,{(\sgi-1)*.055}) -- (-.12,{(\sgi-1)*.055+.03}) arc[start angle=180,end angle=360,x radius=.12,y radius=.045]
        -- (.12,{(\sgi-1)*.055}) arc[start angle=0,end angle=-180,x radius=.12,y radius=.045] -- cycle;
      \filldraw[draw=sggold!70!black,fill=sggold!35,line width=.4pt]
        (0,{(\sgi-1)*.055+.03}) ellipse[x radius=.12,y radius=.045];}
  \end{scope}%
}

%% \sgcard[s]{(x,y)}{colour}: hint/task card with a dot and two lines.
\newcommand{\sgcard}[3][1]{%
  \begin{scope}[shift={#2},scale=#1]
    \filldraw[draw=#3,fill=white,line width=.55pt,rounded corners=1.5pt] (-.15,-.2) rectangle (.15,.2);
    \fill[#3] (-.05,.09) circle (.045);
    \draw[#3!80,line width=.5pt] (-.09,-.03)--(.09,-.03) (-.09,-.11)--(.05,-.11);
  \end{scope}%
}

%% \sgbubble[style]{(x,y)}{(tail tip)}{text}: speech bubble with a tail.
\newcommand{\sgbubble}[4][]{%
  \node[draw=black!45,fill=white,rounded corners=3pt,inner sep=2pt,line width=.45pt,
    sgfont=6.8,align=center,#1] (sgbub) at #2 {#4};%
  \filldraw[draw=black!45,fill=white,line width=.45pt]
    ($(sgbub.south)+(-.07,.03)$) -- #3 -- ($(sgbub.south)+(.07,.03)$);%
  \node[draw=none,fill=white,rounded corners=3pt,inner sep=1.5pt,sgfont=6.8,align=center,#1]
    at (sgbub) {#4};%
}

%% \sgcheck[s]{(x,y)} and \sgcross[s]{(x,y)}: result marks.
\newcommand{\sgcheck}[2][1]{%
  \draw[sggreen,line width=1.6pt,line cap=round,line join=round,shift={#2},scale=#1]
    (-.13,0)--(-.04,-.1)--(.15,.13);}
\newcommand{\sgcross}[2][1]{%
  \draw[sgastra,line width=1.6pt,line cap=round,shift={#2},scale=#1]
    (-.1,-.1)--(.1,.1) (-.1,.1)--(.1,-.1);}

%% \sgdoor[s]{(x,y)}: open door, origin = bottom centre of the frame.
\newcommand{\sgdoor}[2][1]{%
  \begin{scope}[shift={#2},scale=#1]
    \filldraw[draw=black!55,fill=black!6,line width=.6pt] (-.22,0) rectangle (.22,.78);
    \filldraw[draw=black!55,fill=sggold!25,line width=.6pt] (.22,0)--(.42,-.06)--(.42,.84)--(.22,.78)--cycle;
    \fill[black!55] (.37,.38) circle (.025);
  \end{scope}%
}

%% \sgnum{colour}{n}: filled circle with a white number, for step headers.
\newcommand{\sgnum}[2]{\tikz[baseline=(sgn.base)]\node[circle,fill=#1,text=white,inner sep=0pt,
  minimum size=10.5pt,font=\sffamily\bfseries\fontsize{7}{7}\selectfont] (sgn) {#2};}

%% \sgstep{(x,y)}{colour}{n}{title}: numbered step header centred at (x,y).
\newcommand{\sgstep}[4]{%
  \node[inner sep=0pt,font=\sffamily\bfseries\fontsize{8.6}{9}\selectfont,text=#2] at #1
    {\sgnum{#2}{#3}\hspace{3pt}#4};}

%% \sgchev{(x,y)}: grey chevron between panels.
\newcommand{\sgchev}[1]{%
  \draw[black!35,line width=1.3pt,line cap=round,line join=round,shift={#1}]
    (-.06,.11)--(.06,0)--(-.06,-.11);}

%% \sggauge[s]{(x,y)}{colour}{level}: vertical token gauge, bottom centre at (x,y),
%% 1.6 tall; level in 0..1 is the filled share.
\newcommand{\sggauge}[4][1]{%
  \begin{scope}[shift={#2},scale=#1]
    \fill[black!45] (-.06,1.6) rectangle (.06,1.66);
    \fill[#3!70] (-.13,0) rectangle (.13,{1.6*#4});
    \draw[black!55,line width=.55pt,rounded corners=1pt] (-.13,0) rectangle (.13,1.6);
  \end{scope}%
}

%% \sgpack[s]{(x,y)}{colour}: one refill pack (box with a plus), centred at (x,y).
\newcommand{\sgpack}[3][1]{%
  \begin{scope}[shift={#2},scale=#1]
    \filldraw[draw=#3!80!black,fill=#3!18,line width=.6pt,rounded corners=2pt] (-.24,-.18) rectangle (.24,.18);
    \fill[#3!80!black] (-.11,-.025) rectangle (.11,.025) (-.025,-.11) rectangle (.025,.11);
    \draw[#3!80!black,line width=.6pt] (-.08,.18)--(-.08,.24)--(.08,.24)--(.08,.18);
  \end{scope}%
}

%% \sgcut[s]{(x,y)}{i}{cells}: mark cell i of a token bar (same geometry as \sgbar)
%% as just taken: red dashed outline with a slash.
\newcommand{\sgcut}[4][1]{%
  \begin{scope}[shift={#2},scale=#1]
    \pgfmathsetmacro{\sgbarw}{#4*.115-.025}
    \pgfmathsetmacro{\sgbx}{-\sgbarw/2+(#3-1)*.115}
    \draw[sgastra,line width=.45pt,dash pattern=on 1pt off .7pt,fill=sgastra!12] (\sgbx,-.04) rectangle ++(.09,.08);
    \draw[sgastra,line width=.45pt] (\sgbx,-.04)--++(.09,.08);
  \end{scope}%
}

%% \sgnoexit[s]{(x,y)}: door under a red no-entry sign (no EXIT), bottom centre at (x,y).
\newcommand{\sgnoexit}[2][1]{%
  \begin{scope}[shift={#2},scale=#1]
    \filldraw[draw=black!55,fill=black!6,line width=.6pt] (-.2,.05) rectangle (.2,.75);
    \filldraw[draw=black!55,fill=sggold!25,line width=.6pt] (.2,.05)--(.36,0)--(.36,.8)--(.2,.75)--cycle;
    \draw[sgastra,line width=1.3pt] (.05,.4) circle (.44);
    \draw[sgastra,line width=1.3pt] ($(.05,.4)+(135:.44)$)--($(.05,.4)+(-45:.44)$);
  \end{scope}%
}
